\section{Research Participant Information Sheet Template}
\label{app:info-sheet}

\textbf{Introduction}

The purpose of this form is to provide you with information, so you can decide whether to participate in this study. Any questions you may have will be answered by the researcher or by the other contact persons provided below. Once you are familiar with the information on this sheet and have asked any questions you may have, you can decide whether or not to participate. If you agree, you will be asked to fill in the consent form for this study or record your consent verbally.

\begin{longtable}{>{\raggedright\arraybackslash}p{0.28\linewidth} >{\raggedright\arraybackslash}p{0.65\linewidth}}
		\hline
		\textbf{\makecell[lt]{Research \\ title}} & Leveraging Carbon Markets to Scale Up Rural Solar Mini Grids and Accelerate Sustainable Development in Sub-Saharan Africa. \textbf{Alternative Title:} How carbon finance can help fund rural electricity projects in Africa. \\
		\hline
		\textbf{\makecell[lt]{Name and contact \\ details of \\ researcher}} & Name: Flora Aluoch Oloo Email: 684134@soas.ac.uk or floraoloo24@gmail.com \\
		\hline
		\textbf{\makecell[lt]{Name and contact \\ details of Principal \\ Investigator}} & Supervisor: Anabelle De Freece Email: af12@soas.ac.uk \\
		\hline
		\textbf{\makecell[lt]{What type of \\ research project is \\ this?}} & Non-Funded Research Project \\
		\hline
		\textbf{\makecell[lt]{Who is funding \\ this research \\ project?}} & Self-funded Masters Research Project \\
		\hline
		\textbf{\makecell[lt]{Who else is \\ involved with the \\ research project?}} & This research is conducted solely by the researcher at SOAS University of London. No other organizations have access to the data collected, though anonymized findings may be shared with development organizations and policy makers to inform practice. \\
		\hline
		\textbf{\makecell[lt]{What is the \\ research project's \\ purposes?}} & Many rural communities in Sub-Saharan Africa lack access to electricity, which limits economic development and quality of life. At the same time, there are financing challenges that make it difficult to build solar mini-grids in these communities. This research explores whether carbon markets could help finance more rural electricity projects. The study aims to understand the opportunities and barriers for using carbon finance to scale up rural solar mini-grids, assess the potential benefits for communities, and identify what policies would make this approach successful. \\
		\hline
		\textbf{\makecell[lt]{Why have I been \\ chosen?}} & You have been selected because of your expertise and experience in solar mini-grid development and rural electrification in Sub-Saharan Africa. Your insights are valuable for understanding how carbon finance could work in practice for rural electrification projects. Participants in the research projects cut across different stakeholder groups including developers, policy makers, investors, and community members. \\
		\hline
		\textbf{\makecell[lt]{Do I have to take \\ part?}} & Participation in this research is entirely voluntary. You can choose not to answer any questions you're uncomfortable with, or withdraw from the study at any time without giving a reason. If you withdraw, any data you have provided will be deleted unless you specifically agree to its continued use. \\
		\hline
		\textbf{\makecell[lt]{What will happen \\ to me if I take \\ part?}} & You will participate in a 30--45-minute interview (in-person, by phone, or video call) where I will ask about your experiences, perspectives, and opinions on carbon finance and mini-grid development. Questions will cover topics like financing challenges, carbon market mechanisms, policy frameworks, and community impacts. \\
		\hline
		\textbf{\makecell[lt]{Will I be recorded \\ and how will the \\ recordings be \\ used?}} & With your permission, interviews will be audio recorded to ensure accurate data collection. Recordings will only be used for research analysis and will be transcribed by the researcher or a professional transcription service bound by confidentiality agreements. Video recording will only be used if the interview is conducted online and you consent. \\
		\hline
		\textbf{\makecell[lt]{Risks and Benefits \\ of participation}} & \textbf{Benefits:}
		\begin{itemize}[noitemsep,topsep=2pt]
			\item Contributing to research that could help improve energy access in rural Africa
			\item Opportunity to share your expertise and influence policy recommendations
			\item Access to research findings that may inform your work
		\end{itemize}
		\textbf{Risks:}
		\begin{itemize}[noitemsep,topsep=2pt]
			\item Minimal risk beyond those of normal professional conversations
			\item Some questions may touch on business-sensitive topics, but you can decline to answer any question
			\item Time commitment required for participation
		\end{itemize}
		No direct financial benefits are provided for participation. \\
		\hline
		\textbf{\makecell[lt]{What if Something \\ Goes Wrong?}} & If you have concerns about any aspect of this research, please contact me directly using the details above. If you wish to make a formal complaint or have concerns that cannot be resolved through direct contact, you can contact Flora Aluoch Oloo 684134@soas.ac.uk or Anabelle De Freece Email: af12@soas.ac.uk \\
		\hline
		\textbf{\makecell[lt]{Where will \\ information I \\ proved be \\ transferred to?}} & Data will be collected in Kenya and Nigeria for analysis by the researcher at SOAS University of London. All data transfers will comply with international data protection standards. \\
		\hline
		\textbf{\makecell[lt]{How will \\ information I \\ provide be kept \\ secure?}} & Digital recordings and transcripts will be stored on password-protected, encrypted devices and physical documents used during community engagements will be kept in locked filing cabinets and shredded within 1 year of project completion. Only the researcher will have access to identifiable data. \\
		\hline
		\textbf{\makecell[lt]{Will I be kept \\ anonymous in this \\ research project?}} & Yes, you will be anonymized in all research outputs. Steps to protect your anonymity include: \begin{itemize}
			\item Removing names, organizations, and specific locations from transcripts and publications
			\item Using general descriptors (e.g., "mini-grid developer in East Africa" rather than specific company names)
			\item Aggregating data where possible to prevent identification
			\item Reviewing all quotes and examples to ensure they cannot be traced back to you
		\end{itemize} Note: Due to the small size of the mini-grid sector in some regions, complete anonymity cannot be guaranteed if you provide very specific details about unique projects or circumstances. You can indicate if any information should be treated as particularly sensitive. \\
		\hline
		\textbf{\makecell[lt]{What will happen \\ to the results of \\ this research \\ project?}} & Research results will be published in a Masters thesis (available through SOAS library and online repositories as well as academic journal articles). All published materials will maintain participant anonymity. You can request a summary of findings once the research is completed. The research aims to inform policy makers, development organizations, and practitioners working on rural electrification in Sub-Saharan Africa. \\
		\hline
\end{longtable}

\textbf{Data Protection Privacy Notice}

The data controller for this project will be SOAS University of London. The SOAS Data Protection Officer provides oversight of SOAS activities involving the processing of personal data and can be contacted at dataprotection@soas.ac.uk.

Your personal data will be processed for the purposes outlined in this Information Sheet. The legal basis that would be used to process your personal data under UK Data Protection Legislation is the performance of a task in the public interest or in our official authority as a controller. However, for ethical reasons we need your consent to take part in this research project. You can provide your consent for the use of your personal data in this project by completing the consent form that has been provided for you or via audio recording of the information sheet and consent form content.

\textbf{Your Rights}

You have the right to request access under the UK-GDPR to the information which SOAS holds about you. Further information about your rights under the Regulation and how SOAS handles personal data is available on the Data Protection pages of the SOAS website (\url{http://www.soas.ac.uk/infocomp/dpa/index.html}), and by contacting the Information Compliance Manager at the following address: Information Compliance Manager, SOAS, Thornhaugh Street, Russell Square, London WC1H 0XG, United Kingdom (e-mail to: dataprotection@soas.ac.uk).

If you are concerned about how your personal data is being processed, please contact SOAS In the first instance at dataprotection@soas.ac.uk. If you remain unsatisfied, you may wish to contact the Information Commissioner's Office (ICO). Contact details, and details of data subject rights, are available on the ICO website at: \url{https://ico.org.uk/for-organisations/data-protection-reform/overview-of-the-gdpr/individuals-rights/}

\textbf{Copyright Notice}

The consent form asks you to waive copyright so that SOAS and the researcher can edit, quote, disseminate, publish (by whatever means) your contribution to this research project in the manner described to you by the researcher during the consent process. 

\textbf{Contact for Further Information}

Name: Flora Aluoch Oloo Email: 684134@soas.ac.uk or floraoloo24@gmail.com

Thank you for reading this information sheet and for considering taking part in this research study.